% Includable proof fragment. Requires amsmath, amssymb, amsthm, mathtools.
% Environments theorem, lemma, proposition, corollary assumed.
% This is a new derivation relative to the inspected current draft corpus.
\section{Replacing primitive numerical equality by two monus equations}
\label{sec:monus-basis}

We work initially in ordinary intuitionistic first-order logic with actual
numerical equality. The exact numerical language consists of $0,S$, the
staged tied-seed constructor $B$, and a binary symbol $d$. For each fixed
finite earlier-stage body $h(u,v,\vec p)$ the two $B$ schemas are
\[
 B_h(0,\vec p)=S h(0,0,\vec p),\qquad
 B_h(Sn,\vec p)=h(n,B_h(n,\vec p),\vec p).
 \tag{B}
\]
The two additional computation schemas are
\[
 d(x,0)=x,\qquad d(x,Sy)=P(d(x,y)).
 \tag{M}\label{monus:M}
\]
Here $P$ is the fixed $B$-term constructed below; it is not a further
primitive. Numerical coverage and separation are explicit:
\[
 \forall x\,(x=0\lor\exists y\,x=Sy),\qquad 0=S0\to\bot.
 \tag{CS}\label{monus:CS}
\]
The sole induction schema is contextual induction on a literal equation:
\[
 \frac{\Gamma\vdash E(0,\vec p)\quad
       \Gamma,E(z,\vec p)\vdash E(Sz,\vec p)}
      {\Gamma\vdash E(t,\vec p)},
 \tag{EI}\label{monus:EI}
\]
with $z$ fresh for $\Gamma,t,\vec p$. All schemas are universally closed
in their free numerical variables. There are no function variables or
uniform evaluation symbols. Call this ordinary first-order theory $T_{Bd}$.

Coverage is an essential initial resource in the proof below. It must not
be claimed as a consequence of the positive computation equations and
literal-equation induction: the coordinatewise square of the standard
model satisfies those positive equations and their equation-induction
instances, while numerical coverage fails. A native calculus with a sound
numerical case rule may supply this resource instead, but its reconstruction
of ordinary logical rules must be justified independently.

\subsection{Elementary terms available before equality reflection}
Put $1=S0$ and define
\[
 C(n,a)=B(n;u,v.a),\qquad W(n)=B(n;u,v.u),\qquad
 P(n)=W(C(n,n)),\qquad Z(n)=C(n,0).
\]
The computation schemas alone give
\[
 C(0,a)=Sa,\quad C(Sn,a)=a,\quad W(0)=1,\quad W(Sn)=n,
\]
and therefore
\[
 P(0)=0,\quad P(Sn)=n,\qquad Z(0)=1,\quad Z(Sn)=0.
 \tag{PZ}\label{monus:PZ}
\]
Successor is injective, by application of $P$. It never equals zero:
$Sx=0$ would imply $x=P(0)=0$, and hence $S0=0$, contradicting
separation. Coverage and the endpoints of $Z$ now prove
\[
 Z(x)=0\lor Z(x)=1,\qquad Z(x)=1\ \longleftrightarrow\ x=0.
 \tag{Z}\label{monus:Z}
\]
These are proofs in intuitionistic logic using the explicit coverage axiom;
they do not use induction on a disjunction.

Define the one-sided gate
\[
 G(n,a)=P\bigl(P(B(a;u,v.C(Z(n),v)))\bigr).
\]
Literal-equation induction on $a$ proves
\[
 G(0,a)=0,\qquad G(Sn,a)=a.
 \tag{G}\label{monus:G}
\]
For the zero control the inner iterator has initial value $1$ and identity
update, so it is always $1$. For a successor control its initial value is
$2$ and its update is successor, so its value at $a$ is $S(Sa)$. Applying
$P$ twice gives the two equations. In particular, $G(1,a)=a$.

The inherited zero-start recursion macro is also available without
numerical equality:
\[
 L_h(n,\vec p)=B\bigl(Sn;r,s.G(r,h(Pr,s,\vec p))\bigr).
\]
Its first actual update is zero, and the subsequent updates are $h$;
therefore its equations are
\[
 L_h(0,\vec p)=0,\qquad
 L_h(Sn,\vec p)=h(n,L_h(n,\vec p),\vec p).
 \tag{L}
\]
No arbitrary initialization has been assumed here.

\subsection{Four preliminary monus identities}

\begin{lemma}\label{monus:preliminary}
The theory proves, using only literal-equation induction,
\begin{align}
 d(0,y)&=0,\label{monus:leftzero}\\
 P(d(Sx,y))&=d(x,y),\label{monus:commute}\\
 d(Sx,Sy)&=d(x,y),\label{monus:diagonal}\\
 d(x,x)&=0,\qquad d(Sx,x)=1.\label{monus:diagonal-values}
\end{align}
\end{lemma}
\begin{proof}
For \eqref{monus:leftzero}, induct on $y$ using $P0=0$.
For \eqref{monus:commute}, keep $x$ fixed and induct on $y$.
The base is $P(Sx)=x$. At the step,
\[
 P(d(Sx,Sy))=P(P(d(Sx,y)))
           =P(d(x,y))=d(x,Sy),
\]
where the middle equality applies $P$ to the induction equation.
Equation \eqref{monus:diagonal} follows immediately from
\eqref{monus:M} and \eqref{monus:commute}.
Finally induct on $x$ in the two literal equations of
\eqref{monus:diagonal-values}. Their bases are $d(0,0)=0$ and
$d(1,0)=1$, and their steps are exactly \eqref{monus:diagonal}.
No induction hypothesis is generalized or specialized to a changed
passive parameter.
\end{proof}

\subsection{The gated complement identity}
The key step is an equation invariant that permits the ordinary
conditional complement argument before quantifier-free induction is
available.

\begin{lemma}[Gated complement]\label{monus:complement}
For all $x,y$ the theory proves
\[
 G\bigl(Z(d(y,x)),d(x,d(x,y))\bigr)
       =G\bigl(Z(d(y,x)),y\bigr).
 \tag{GC}\label{monus:GC}
\]
Consequently,
\[
 d(y,x)=0\quad\longrightarrow\quad d(x,d(x,y))=y.
 \tag{C}\label{monus:C}
\]
\end{lemma}
\begin{proof}
Keep $x$ fixed and apply literal-equation induction to
\eqref{monus:GC} as an equation in $y$. At $y=0$, the gate control
is $Z(d(0,x))=1$, and both output arguments are zero by
$d(x,0)=x$ and $d(x,x)=0$.

For the step, retain the induction equation and set
$a=d(Sy,x)$. Use coverage on $a$. If $a$ is a successor, $Z(a)=0$
and both sides of the new equation are zero by \eqref{monus:G}.
It remains to consider $a=0$.

By \eqref{monus:commute},
\[
 d(y,x)=P(d(Sy,x))=0.
\]
Thus the induction equation, with both gates simplified at control $1$,
gives
\[
 d(x,t)=y,\qquad t=d(x,y).
 \tag{IH}\label{monus:IH}
\]
If $t=0$, then $d(x,t)=x$, so $x=y$. But this makes
$a=d(Sy,y)=1$, contradicting $a=0$. Hence $t\ne0$.
Coverage therefore provides $u$ with $t=Su$. The monus successor
schema identifies
\[
 u=P(t)=d(x,Sy).
\]
Put $v=d(x,u)$. Equation \eqref{monus:IH} and the successor schema give
\[
 P(v)=P(d(x,u))=d(x,Su)=d(x,t)=y.
\]
If $v=0$, then $y=0$. Hence $t=d(x,0)=x=Su$, and
\[
 v=d(x,u)=d(Su,u)=1,
\]
a contradiction. Thus $v\ne0$, and coverage gives $v=Sz$.
Since $P(v)=y$, the predecessor equations give $z=y$, so $v=Sy$.
Together with $u=d(x,Sy)$ this is
\[
 d(x,d(x,Sy))=Sy,
\]
which is exactly the required step with gate control $Z(a)=1$.
All contradictions use only $0\ne1$ and intuitionistic falsity elimination.
All case splits occur inside the step derivation with the original
induction equation retained. The only induction matrix used was the
single literal equation \eqref{monus:GC}.

Finally if $d(y,x)=0$, the control of both gates in
\eqref{monus:GC} is $1$, giving \eqref{monus:C}.
\end{proof}

\begin{corollary}[Antisymmetry before order theory]\label{monus:antisymmetry}
The theory proves
\[
 d(x,y)=0\ \wedge\ d(y,x)=0\quad\longrightarrow\quad x=y.
\]
\end{corollary}
\begin{proof}
Equation \eqref{monus:C} gives $d(x,d(x,y))=y$. Substituting
$d(x,y)=0$ and then $d(x,0)=x$ gives $x=y$.
\end{proof}

\subsection{A characteristic of actual equality}
Define the fixed numerical term
\[
 \epsilon(x,y)=G\bigl(Z(d(x,y)),Z(d(y,x))\bigr).
 \tag{E}\label{monus:E}
\]

\begin{theorem}[Whole-domain equality reflection]\label{monus:reflection}
The theory $T_{Bd}$ proves
\[
 \epsilon(x,y)=0\lor\epsilon(x,y)=1,
 \qquad
 \epsilon(x,y)=1\ \longleftrightarrow\ x=y.
\]
These statements concern actual equality throughout the domain.
\end{theorem}
\begin{proof}
Each input of the gate is zero or one by \eqref{monus:Z}. The two
gate equations prove the displayed range. They also prove
\[
 \epsilon(x,y)=1\quad\longleftrightarrow\quad
                  d(x,y)=0\wedge d(y,x)=0:
\]
if the first gate input were zero, the output would be zero; when it
is one, the output is the second input. Use $0\ne1$ and
\eqref{monus:Z} in both directions. Antisymmetry proves the right-to-left
numerical identification required for reflection. Conversely $x=y$
implies both differences are zero by $d(x,x)=0$ and equality substitution,
so the gate output is one.
\end{proof}

In particular numerical equality is decidable, but excluded middle for
arbitrary quantified formulas has not been postulated. Every fixed
quantifier-free formula has a Boolean characteristic term: use
\eqref{monus:E} for atoms, zero for falsity, $Z$ for negation, and $G$
for conjunction; disjunction and implication follow from the Boolean
identities. Implication is legitimate intuitionistically because the
antecedent is decidable. Literal-equation induction on the characteristic
being one yields contextual quantifier-free induction.

\subsection{Consequences for the existing bounded compiler}

The equality term just proved satisfies precisely the range and actual
reflection axioms of the inspected $B+\epsilon$ bounded-initializer
presentation. Its finite descriptions and literal-equation induction
instances translate by substitution of \eqref{monus:E}. Therefore the
existing bounded-initialization construction, cap-coherence argument,
and finite quantifier-free function graphs apply to $T_{Bd}$.
They give a whole-domain interpretation of iPRA with actual equality.

In the reverse direction, iPRA interprets each fixed $B_h$ by primitive
recursion with base $S h(0,0,\vec p)$ and step $h$, and interprets $d$
by ordinary truncated subtraction. It proves the two monus equations,
coverage, separation, and every source equation-induction instance.
The common numerical domain is unchanged.

For the source-side composite, first identify the $d$-free constant-body
$C$ and counter-body $W$ recursors by recurrence uniqueness, and hence
their predecessor macro $P$. The imported primitive-recursive monus
graph is provably equal to the original $d$: induct on its second
argument using its defining equations and \eqref{monus:M}.
When a graph-defined function is used in this comparison, expand its
finite quantifier-free graph cases; the quantifier-free induction already
derived above licenses this definitional-extension argument. For each
fixed staged $B$ symbol, recurrence uniqueness similarly identifies the
composite with the original symbol, proceeding by external induction on
the finite description. The iPRA-side composite follows from primitive-
recursion uniqueness in the same way. Thus the inherited compiler proof
extends to full-domain bi-interpretability and equiconsistency.

This is a new arithmetic presentation with a transparent reflection
bootstrap. It is not, by itself, a proof-checker size reduction. Its
ordinary logical rules, coverage, separation, all four computation
schemas, and equation induction must be included in any honest
implementation comparison. A checker for a different inherited calculus
cannot be assigned to this new theory merely because both are
bi-interpretable with iPRA.
